\section{Introducción: repaso del curso y propósito de esta clase}

Esta clase es la última de las seis clases fundamentales de la serie MIT 6.S191 Introduction to Deep Learning. En las cinco clases anteriores, el curso presentó de manera sistemática la metodología central del aprendizaje profundo: desde las redes neuronales completamente conectadas básicas, el modelado de secuencias (RNN / Transformer), las redes neuronales convolucionales (CNN), hasta los modelos generativos (VAE / GAN) y el aprendizaje por refuerzo profundo. Esta clase, apoyándose en estos fundamentos, se concentrará en dos temas centrales:

\begin{enumerate}
    \item \textbf{las limitaciones del aprendizaje profundo} (Limitations): el dilema de la generalización, los ataques adversarios, la representación de la incertidumbre, entre otros;
    \item \textbf{las nuevas fronteras} (New Frontiers): los avances más recientes representados por Diffusion Models y Large Language Models (LLMs).
\end{enumerate}

\begin{figure}[H]
\centering
\includegraphics[width=0.85\textwidth]{slides-images/slide-01.jpg}
\caption{Tema de esta clase: Deep Learning Limitations and New Frontiers\protect\footnotemark}
\end{figure}
\footnotetext{Slides, página 1.}

Como señaló Amini en la clase, el aprendizaje profundo ya ha traído avances revolucionarios en numerosos campos, como la conducción autónoma, las imágenes médicas, el aprendizaje por refuerzo, el modelado generativo, el procesamiento de lenguaje natural, las finanzas y la seguridad. Pero como profesionales de la tecnología, no basta con entender el poder de estos algoritmos; también debemos reconocer con claridad sus limitaciones, para poder desplegar la IA de manera responsable en el mundo real.

\begin{figure}[H]
\centering
\includegraphics[width=0.85\textwidth]{slides-images/slide-10.jpg}
\caption{Auge del aprendizaje profundo en diversos campos\protect\footnotemark}
\end{figure}
\footnotetext{Slides, página 10.}

\subsection{Resumen de la sección}

Esta clase, como cierre de las clases fundamentales, conecta lo anterior con lo que sigue: repasa la esencia del aprendizaje profundo como aproximador de funciones, señala sus limitaciones y presenta las nuevas fronteras representadas por Diffusion Models y LLM.

%% =====================================================
